\begin{abstract}
\noindent 由监测病例数得到的再生数估计，包括由每周住院数计算的表观更新型比值，通常被认为会随报告病例增多而更加精确。传播强度却随天气、学校校历和人群行为在时期之间变化。本文基于传播存在环境变异、后代服从负二项分布、报告服从二项分布的分支过程推导得出，每期报告 $m$ 例时，平均再生数的单期误差方差为 $V(m)=v_R+A/m$。因此，病例量超过平衡规模 $\mx=A/v_R$ 后精度收益递减；固定预算的最优用法是每期取 $m^*=\sqrt{(c_0/c_1)\mx}$ 例，其中 $c_0/c_1$ 为单个时期与单个病例的成本之比。我们分析了 2020--2026 年美国流感住院数据的 10506 个州周。在四个流行季的早期增长窗口中，更新型比值的相对损失随住院数增加而趋于正截距，按该比值的尺度对应每周 37.4 例住院的平衡规模（95\% CI：19.1, 71.0）。经由数据匹配模拟的两步反演，更新型比值与加权增长比分别对应 0.040（95\% CI：0.026, 0.056）和 0.043（95\% CI：0.025, 0.064）的潜伏环境方差，两者相互一致（比值 0.93）。两个估计都以 3.6 天的代际间隔和相互独立的周冲击为条件，数据的滞后一阶自相关支持后一条件。截距的 41\%--51\% 为同一日历周内各州所共有。强制报告下，机构报告完整度使截距变化 2\%（上界 5.3\%；参考报告波动取 0.10 时，等效界为截距的 17\%）。2024 年自愿报告期，周增长的离散程度是 2023 年同期的 1.58 倍（95\% CI：0.95, 2.23）。在每周住院数超过平衡规模的情况下（增长窗口中占 55\% 的州周），增加时期数与区域数比增加每期病例更能提高精度。较高的机构参与度与较低的波动相伴。

\medskip
\noindent\textbf{关键词：}再生数；环境随机性；监测设计；报告异质性；流感住院；收益递减
\end{abstract}
